\section*{Глава 2. Проектирование веб-приложения}
\addcontentsline{toc}{section}{Глава 2. Проектирование веб-приложения}
\stepcounter{section}

\subsection{Анализ предметной области и целевой аудитории}

Разрабатываемая система \textit{Table \& Time} предназначена для автоматизации
процессов бронирования столиков и оформления заказов в ресторане.
Основной задачей системы является объединение функций цифрового меню,
онлайн-бронирования и внутренней панели управления персонала в рамках
единого веб-приложения.

Современные заведения общественного питания нуждаются в цифровых
инструментах, позволяющих упростить взаимодействие между клиентами
и персоналом ресторана. Использование веб-приложений позволяет
сократить время ожидания обслуживания, снизить нагрузку на сотрудников
и повысить удобство оформления заказов и бронирований.

Разрабатываемая система ориентирована как на посетителей ресторана,
использующих приложение для выбора блюд и бронирования столиков,
так и на сотрудников заведения, осуществляющих обработку заказов
и управление внутренними процессами ресторана.

Целевой аудиторией системы являются посетители ресторана, а также сотрудники
заведения. В системе выделяются четыре основные пользовательские роли:

\begin{itemize}
	\item \textbf{Гость (Guest)} --- неавторизованный пользователь, имеющий
	      доступ к просмотру меню, информации о ресторане и схемы зала;
	\item \textbf{Пользователь (User)} --- зарегистрированный клиент,
	      способный бронировать столики, оформлять заказы и просматривать историю
	      посещений;
	\item \textbf{Официант (Waiter)} --- сотрудник ресторана, управляющий
	      заказами и подтверждением бронирований;
	\item \textbf{Администратор (Admin)} --- пользователь с расширенными
	      полномочиями, осуществляющий управление меню, столами, сотрудниками
	      и параметрами системы.
\end{itemize}

Разделение пользователей по ролям реализовано посредством механизма
RBAC (Role-Based Access Control), обеспечивающего разграничение
доступа к функциональным возможностям приложения.

\subsection{Описание функциональности системы}

Функциональность приложения определяется через набор пользовательских
сценариев взаимодействия с системой.

Основной пользовательский сценарий можно сформулировать следующим образом:

\begin{quote}
	«Я, как пользователь ресторана, хочу выбрать свободный столик,
	забронировать его на определённое время и оформить заказ через
	веб-приложение, чтобы сократить время ожидания в заведении».
\end{quote}

Система поддерживает следующие функциональные возможности:

\begin{itemize}
	\item просмотр интерактивного меню ресторана;
	\item фильтрацию блюд по категориям;
	\item мультиязычный интерфейс (RU/EN);
	\item бронирование столиков через интерактивную карту зала;
	\item оформление заказов различных типов:
	      \begin{itemize}
		      \item заказ в зале (\texttt{DINE\_IN});
		      \item доставка (\texttt{DELIVERY});
		      \item самовывоз (\texttt{TAKEAWAY});
	      \end{itemize}
	\item управление статусами заказов персоналом;
	\item административное управление меню, пользователями и столами;
	\item хранение истории заказов и бронирований.
\end{itemize}

На рисунке~\ref{fig:use_case} представлена диаграмма вариантов
использования системы.

\begin{figure}[H]
	\centering
	\includegraphics[width=0.7\textwidth]{use_case_diagram}
	\caption{Диаграмма вариантов использования системы}
	\label{fig:use_case}
\end{figure}

\subsection{Архитектура веб-приложения}

Разрабатываемая система построена на основе клиент-серверной архитектуры.
Frontend-часть реализована с использованием фреймворка Next.js 15
(App Router), а backend --- на базе NestJS.

Архитектура приложения включает следующие уровни:

\begin{itemize}
	\item \textbf{Frontend} --- пользовательский интерфейс, реализованный
	      на React и TypeScript;
	\item \textbf{Backend API} --- REST API, реализованный на NestJS;
	\item \textbf{База данных} --- PostgreSQL;
	\item \textbf{ORM-уровень} --- Prisma ORM;
	\item \textbf{Система авторизации} --- JWT + Refresh Token.
\end{itemize}

Frontend и backend взаимодействуют через REST API с использованием
формата JSON.

На рисунке~\ref{fig:architecture} представлена общая архитектура
веб-приложения.

\begin{figure}[H]
	\centering
	\includegraphics[width=0.85\textwidth]{architecture_diagram}
	\caption{Архитектура веб-приложения}
	\label{fig:architecture}
\end{figure}

\subsection{Проектирование базы данных}

Для хранения данных была выбрана реляционная СУБД PostgreSQL,
обеспечивающая поддержку транзакций, индексации и связей между сущностями.
Выбор PostgreSQL обусловлен её высокой производительностью, надёжностью
и широкими возможностями по работе со сложными структурами данных.

Использование реляционной модели позволяет обеспечить целостность данных,
а также корректную обработку взаимосвязей между пользователями,
заказами, бронированиями и другими сущностями системы.
Дополнительно PostgreSQL предоставляет механизмы оптимизации запросов,
что особенно важно при обработке операций бронирования столиков
и поиска активных заказов.

В качестве ORM используется Prisma, предоставляющая типобезопасное
взаимодействие с базой данных. Prisma автоматически генерирует клиент
для работы с моделями данных, что снижает вероятность ошибок при
выполнении запросов и упрощает разработку backend-части приложения.

Использование Prisma также позволяет централизованно описывать структуру
базы данных через схему данных, поддерживать миграции и обеспечивать
синхронизацию между приложением и СУБД.

Основными сущностями системы являются:

\begin{itemize}
	\item пользователи (\texttt{users});
	\item блюда (\texttt{dishes});
	\item столики (\texttt{tables});
	\item заказы (\texttt{orders});
	\item бронирования (\texttt{reservations});
	\item позиции заказа (\texttt{order\_items});
	\item рестораны (\texttt{restaurants});
	\item refresh-токены (\texttt{refresh\_tokens}).
\end{itemize}

Связи между сущностями представлены на ER-диаграмме
(рисунок~\ref{fig:er_diagram}).

\begin{figure}[H]
	\centering
	\includegraphics[width=1\textwidth]{er_diagram}
	\caption{ER-диаграмма базы данных}
	\label{fig:er_diagram}
\end{figure}

Структура таблицы заказов представлена в
таблице~\ref{tab:orders}.

\begin{longtable}{|l|l|p{3.5cm}|p{5.5cm}|}
	\caption{Структура таблицы заказов (\texttt{orders})}
	\label{tab:orders}                                                                                                      \\

	\hline
	\textbf{Поле}      & \textbf{Тип} & \textbf{Ограничение} & \textbf{Описание}                    \\ \hline
	\endfirsthead

	\hline
	\textbf{Поле}      & \textbf{Тип} & \textbf{Ограничение} & \textbf{Описание}                    \\ \hline
	\endhead

	id                  & UUID         & PRIMARY KEY          & Уникальный идентификатор заказа      \\ \hline
	user\_id            & UUID         & FOREIGN KEY          & Пользователь, создавший заказ        \\ \hline
	type                & ENUM         & NOT NULL             & Тип заказа                           \\ \hline
	status              & ENUM         & DEFAULT CREATED      & Текущий статус заказа                \\ \hline
	table\_id           & UUID         & NULL                 & Столик для заказа в зале             \\ \hline
	delivery\_address   & TEXT         & NULL                 & Адрес доставки                       \\ \hline
	reservation\_id     & UUID         & NULL                 & Связанное бронирование               \\ \hline
	payment\_method     & ENUM         & NULL                 & Метод оплаты                         \\ \hline
	total\_amount       & DECIMAL      & NOT NULL             & Общая сумма заказа                   \\ \hline
	scheduled\_at       & TIMESTAMP    & NULL                 & Время самовывоза                     \\ \hline
	created\_at         & TIMESTAMP    & DEFAULT NOW()        & Время создания                       \\ \hline
	updated\_at         & TIMESTAMP    & DEFAULT NOW()        & Время изменения                      \\ \hline
\end{longtable}

Для повышения производительности системы были реализованы
индексы по наиболее часто используемым полям:

\begin{itemize}
	\item \texttt{[status, scheduledAt]};
	\item \texttt{[userId]};
	\item \texttt{[tableId]};
	\item \texttt{[restaurantId, startTime]}.
\end{itemize}

\subsection{Модель обработки бронирований}

Ключевой особенностью системы является механизм проверки
пересечения временных интервалов бронирований.

При создании брони выполняется проверка:

\begin{equation}
	startTime < requestedEndTime
	\land
	endTime > requestedStartTime.
\end{equation}

Данная проверка позволяет исключить возможность
одновременного бронирования одного столика несколькими пользователями.

Для ускорения поиска конфликтующих записей используется
составной индекс:

\begin{equation}
	(tableId,\ startTime,\ endTime).
\end{equation}

\subsection{Проектирование системы авторизации}

Для обеспечения безопасности приложения используется
JWT-аутенти- фикация с применением Access и Refresh Token.

Система авторизации включает:

\begin{itemize}
	\item access token со временем жизни 15 минут;
	\item refresh token со временем жизни 7 дней;
	\item хранение refresh token в базе данных;
	\item защиту маршрутов через Guards;
	\item разграничение доступа по ролям.
\end{itemize}

Авторизация реализована посредством следующих компонентов:

\begin{itemize}
	\item \texttt{JwtAuthGuard};
	\item \texttt{RolesGuard};
	\item \texttt{OrderAccessGuard};
	\item стратегий \texttt{JwtStrategy} и \texttt{RefreshStrategy}.
\end{itemize}

\subsection{Проектирование пользовательского интерфейса}

Пользовательский интерфейс системы реализован с применением
концепции glassmorphism-дизайна и адаптивной вёрстки.

В качестве основных технологий интерфейса используются:

\begin{itemize}
	\item React;
	\item Next.js App Router;
	\item Tailwind CSS;
	\item TypeScript.
\end{itemize}

Для формирования единого визуального стиля были определены
фирменные цвета приложения:

\begin{itemize}
	\item amber --- \#FFA239;
	\item cream --- \#FCF9EA;
	\item sage --- \#97A87A;
	\item espresso --- \#4B2E2B.
\end{itemize}

Интерфейс поддерживает адаптацию под мобильные устройства,
анимации пользовательских элементов и мультиязычный режим.

\subsection{Разработка макетов страниц}

На этапе проектирования были разработаны макеты
ключевых страниц приложения.

\subsubsection{Главная страница}

Главная страница выполняет роль лендинга и содержит:

\begin{itemize}
	\item hero-секцию;
	\item информацию о ресторане;
	\item галерею изображений;
	\item CTA-блок бронирования;
	\item навигационное меню.
\end{itemize}

Визуальный стиль страницы построен на использовании
glassmorphism-эффектов, полупрозрачных карточек и
градиентных фонов.

Макет главной страницы представлен на
рисунке~\ref{fig:landing_mockup}.
\begin{figure}[H]
	\centering
	\includegraphics[width=0.55\textwidth]{landing_mockup}
	\caption{Макет главной страницы}
	\label{fig:landing_mockup}
\end{figure}

\subsubsection{Страница меню}

Страница меню предназначена для просмотра блюд,
фильтрации по категориям и формирования корзины заказа.

Интерфейс страницы включает:

\begin{itemize}
	\item строку поиска;
	\item фильтрацию по категориям;
	\item карточки блюд;
	\item интерактивную корзину заказа;
	\item адаптивную сетку отображения.
\end{itemize}

Макет страницы меню представлен на
рисунке~\ref{fig:menu_mockup}.

\begin{figure}[H]
	\centering
	\includegraphics[width=0.9\textwidth]{menu_mockup}
	\caption{Макет страницы меню}
	\label{fig:menu_mockup}
\end{figure}

\subsubsection{Страница бронирования}

Страница бронирования реализует интерактивную
схему зала ресторана.

Основными элементами интерфейса являются:

\begin{itemize}
	\item интерактивная карта столиков;
	\item выбор даты и времени;
	\item фильтрация по количеству гостей;
	\item визуальное отображение статусов столов;
	\item модальное окно бронирования.
\end{itemize}

Статусы столиков отображаются цветовой индикацией:

\begin{itemize}
	\item зелёный --- свободен;
	\item жёлтый --- забронирован;
	\item красный --- занят;
	\item серый --- уборка.
\end{itemize}

Макет страницы бронирования представлен на
рисунке~\ref{fig:reservation_mockup}.

\begin{figure}[H]
	\centering
	\includegraphics[width=0.95\textwidth]{reservation_mockup}
	\caption{Макет страницы бронирования}
	\label{fig:reservation_mockup}
\end{figure}

\subsection{Вывод по главе}

В результате проектирования была разработана архитектура
веб-приложе- ния \textit{Table \& Time}, включающая
модульный backend на NestJS, frontend на Next.js,
систему ролевого доступа и реляционную модель данных.

Были определены основные пользовательские сценарии,
спроектирована структура базы данных, реализована
система бронирования и разработаны макеты ключевых
страниц приложения.

Полученные результаты создают основу для последующей
реализации и тестирования программной системы.